\section[Rédaction Bourbaki \redac{596} : algèbres, produits croisés, modules induits]{Rédaction Bourbaki \redac{596} : algèbres, produits croisés, modules induits\protect\verifie{Redaction596Algebras}}\label{sec:bourbaki-596-algebres}

La section~\ref{sec:bourbaki-596-hypotheses} a relevé dans la rédaction \redac{596} trois énoncés auxquels manque une hypothèse. La lecture complète des §§~7, 9 et~12 en donne six autres, de deux sortes. Quatre énoncés sont faux tels qu'ils sont écrits : la remarque~3 du §~9, n\textsuperscript{o}\,4 ; l'assertion~(b) de la proposition~6 du §~12 ; une glose au crayon du corollaire~4 du §~7 ; le corollaire~1 de la proposition~26 du §~12 quand le corps n'est pas parfait. Deux autres écarts sont des fautes de calcul dans des démonstrations dont l'énoncé reste juste : la proposition~12\,a) du §~9 et la proposition~3 du §~12. La transcription reste locale ; on n'en cite aucun passage.

\subsection*{Énoncés faux tels qu'ils sont écrits}

\emph{§~9, n\textsuperscript{o}\,4, remarque~3 (p.~147 du fac-similé).} Sur un corps valué localement compact, toute algèbre simple centrale serait isomorphe à un produit croisé galoisien cyclique $A[\mathcal E, L]$. Un tel produit croisé est de dimension $[L:K]^2$ et contient $L$ comme sous-corps commutatif. La remarque~2 qui précède exclut déjà les corps algébriquement clos ; il manque l'hypothèse que la valeur absolue soit ultramétrique.

\emph{§~12, n\textsuperscript{o}\,3, proposition~6\,(b) (p.~261).} Soient $H$ un sous-groupe de $G$ et $F$ un $H$-module. Le $k$-module $\mathrm{Ind}(F)$ serait libre si et seulement si $F$ l'est. La décomposition~(40) montre seulement que $\mathrm{Ind}(F)$ est, comme $k$-module, somme directe de $[G:H]$ copies de $F$ ; cela suffit pour l'assertion~(a) sur les modules projectifs, non pour la liberté.

\emph{§~7, n\textsuperscript{o}\,3, corollaire~4, glose au crayon (p.~104).} Pour $S_0$ et $S_1$ simples, le corollaire tapé dit que $S = S_0 \otimes S_1$ est simple si et seulement si $D = D_0 \otimes D_1$ est un corps. Une main au crayon, non identifiée, traduit ce dernier cas par la condition que le centre $Z = Z_0 \otimes Z_1$ de $D$ soit un corps. L'énoncé tapé est juste ; la glose ne l'est pas.

\emph{§~12, n\textsuperscript{o}\,11, corollaire~1 de la proposition~26 (p.~286).} Avec $\Gamma = \mathrm{Gal}(\bar k/k)$, $\bar k$ une clôture algébrique de $k$, le corollaire identifie les invariants de $\Gamma \times \{e\}$ dans $\bar k \otimes R_{\bar k}(G)$ à $k \otimes R_{\bar k}(G)$. Il faut pour cela que le corps des invariants de $\Gamma$ dans $\bar k$ soit $k$ ; quand $k$ est imparfait, c'est la clôture parfaite de $k$. Il faut prendre une clôture séparable.

\begin{encadre}
\begin{contrexemple}[aux quatre énoncés de la rédaction \redac{596}, tels qu'ils sont écrits]
\begin{enumerate}
\item[a)] Le corps $\mathbf C$ est valué et localement compact. L'algèbre $M_2(\mathbf C)$ est simple centrale, mais n'est pas un produit croisé : il faudrait un corps $L \supset \mathbf C$ avec $[L:\mathbf C] = 2$.
\item[b)] Soient $k = \mathbf Z[\sqrt{-5}]$, $G = \mathbf Z/2$, $H = \{e\}$ et $F = I = (2, 1 + \sqrt{-5})$. Alors $\mathrm{Ind}(I) \cong I \oplus I \cong k^2$ est libre, et $I$ ne l'est pas.
\item[c)] Soient $K = \mathbf R$ et $A_0 = A_1 = S_0 = S_1 = \mathbf H$. Le centre $Z = \mathbf R \otimes_{\mathbf R} \mathbf R = \mathbf R$ est un corps, mais $\mathbf H \otimes_{\mathbf R} \mathbf H$ a des diviseurs de zéro et n'est pas un module simple sur lui-même.
\item[d)] Soient $p$ premier et $k = \mathbf F_p(t)$. Dans une clôture algébrique $\bar k$, une racine $p$-ième $x$ de $t$ est fixée par tout $k$-automorphisme de $\bar k$, et $x \notin k$.
\end{enumerate}
\end{contrexemple}
\end{encadre}

\begin{proof}
a) Une extension finie d'un corps algébriquement clos est triviale : si $[L:\mathbf C] = 2$, $L$ est entier sur $\mathbf C$, donc $\mathbf C \to L$ est bijectif et $[L:\mathbf C] = 1$.

b) Posons $\beta = 1 + \sqrt{-5}$ et $\beta' = 1 - \sqrt{-5} = 2 - \beta$, qui sont dans $I$, avec $\beta\beta' = 6$. Les produits $4$, $2\beta$, $\beta^2 = 2(-2 + \sqrt{-5})$ sont divisibles par~$2$, donc $I \cdot I \subset 2k$. Soient
\[ T = \begin{pmatrix} \beta' & -2 \\ -2 & \beta \end{pmatrix}, \qquad S = \begin{pmatrix} \beta & 2 \\ 2 & \beta' \end{pmatrix}, \qquad ST = TS = 2. \]
L'application $T : k^2 \to I^2$ est injective, car $ST = 2$ et $2$ n'est pas diviseur de zéro ; elle est surjective, car pour $(a, b) \in I^2$ le vecteur $S(a, b)$ est dans $(I\cdot I)^2 \subset 2k^2$, soit $S(a,b) = 2v$, et alors $2\,Tv = TS(a,b) = 2(a,b)$. Donc $I \oplus I \cong k^2$. Si $I$ était libre, il serait de rang au plus~$1$, donc principal, engendré par un $g$ ; la norme $N(g) = a^2 + 5b^2$ diviserait $N(2) = 4$ et $N(\beta) = 6$, donc~$2$. L'équation $a^2 + 5b^2 = 2$ n'a pas de solution entière, et $N(g) = 1$ ferait de $g$ une unité, d'où $I = k$ ; or l'homomorphisme $k \to \mathbf F_2$, $a + b\sqrt{-5} \mapsto a + b$, annule $I$ et pas~$1$.

c) Dans $\mathbf H \otimes \mathbf H$, $(i \otimes 1 - 1 \otimes i)(i \otimes 1 + 1 \otimes i) = i^2 \otimes 1 - 1 \otimes i^2 = 0$. La forme linéaire $x \otimes y \mapsto x_1 y_0$ (coefficient de $i$ dans $x$ fois partie réelle de $y$) vaut $1$ sur les deux facteurs, qui sont donc non nuls. Un anneau simple comme module sur lui-même est un corps, et n'a pas de diviseurs de zéro. Ici $D_i = \mathrm{End}_{\mathbf H}(\mathbf H) \cong \mathbf H^{\circ} \cong \mathbf H$ et $Z_i = \mathbf R$.

d) Pour tout $k$-automorphisme $\sigma$, $\sigma(x)^p = \sigma(t) = t = x^p$, et le Frobenius est injectif, donc $\sigma(x) = x$. Si $x = y \in k$, alors $y^p = t$ ; le degré de $t$ vaut~$1$ et celui de $y^p$ vaut $p\deg y$, ce qui est impossible.
\end{proof}

\noindent Pour a), $\mathbf R$ est un autre contre-exemple ($M_3(\mathbf R)$ demanderait une extension cubique de $\mathbf R$). Sur un corps local non archimédien ou un corps de nombres, l'énoncé est le théorème classique de Hasse--Brauer--Noether.

\subsection*{Fautes dans des démonstrations, énoncés justes}

\emph{§~9, n\textsuperscript{o}\,9, démonstration de la proposition~12\,a) (p.~160).} En caractéristique $p$, l'algèbre $A[a,b)$ est engendrée par $x$, $y$ avec $x^p - x = a$, $y^p = b$, $yx = (x+1)y$, et de même $A[a', b)$ par $u$, $v$. Pour construire un homomorphisme $A[a + a', b) \otimes A[a', 1) \to A[a,b) \otimes A[a',b)$, la démonstration pose $u_1 = 1 \otimes u$ et $v_1 = y \otimes v^{-1}$ et affirme $v_1 u_1 = (u_1 + 1) v_1$. Or $vu = (u+1)v$ entraîne $v^{-1}u = (u - 1)v^{-1}$, d'où $v_1 u_1 = (u_1 - 1)v_1$. Le choix $v_1 = y^{-1} \otimes v$ donne la bonne relation, vérifie $v_1^p = b^{-1} \otimes b = 1$ et commute encore à $x_1 = x \otimes 1 + 1 \otimes u$ et à $y_1 = y \otimes 1$ ; la proposition subsiste.

\emph{§~12, n\textsuperscript{o}\,2, démonstration de la proposition~3 (p.~257).} La démonstration exprime $\psi^p \chi_E$ comme polynôme à coefficients entiers en les $\chi_{\Lambda^q E}$ pour $0 \le q \le p-1$. La formule~(27) fait intervenir $q \le p$ ; déjà $\psi^2\chi = \chi^2 - 2\chi_{\Lambda^2 E}$. Avec la borne $q \le p$ la démonstration est juste.

\begin{encadre}
\begin{contrexemple}[aux deux démonstrations, telles qu'elles sont écrites]
\begin{enumerate}
\item[e)] Sur $\mathbf F_3$, avec $a = a' = 0$ et $b = 1$, soient $x = \operatorname{diag}(0, 1, 2)$ et $y$ la permutation circulaire $e_j \mapsto e_{j-1}$ ; ils vérifient $x^3 = x$, $y^3 = 1$, $yx = (x+1)y$. Avec $u = x$ et $v = y$, les produits de Kronecker $9 \times 9$ donnent $v_1 u_1 = (u_1 - 1)v_1 \neq (u_1 + 1)v_1$ pour $v_1 = y \otimes y^2$, et $v_1 u_1 = (u_1 + 1)v_1$, $v_1 x_1 = x_1 v_1$, $v_1^3 = 1$ pour $v_1 = y^2 \otimes y$.
\item[f)] Soit $G = \mathbf Z/2 = \{1, s\}$. La représentation de permutation sur $\mathbf Q^2$ et la représentation nulle ont en $s$ les mêmes valeurs de $\chi_{\Lambda^0 E} = 1$ et de $\chi_E$ (à savoir $1$ et $0$), mais $(\psi^2\chi_E)(s) = \chi_E(1)$ vaut $2$ pour l'une et $0$ pour l'autre. Aucune fonction de ces deux valeurs, polynomiale ou non, ne donne $\psi^2\chi_E$.
\end{enumerate}
\end{contrexemple}
\end{encadre}

\begin{proof}
e) Le calcul général vaut dans tout produit tensoriel $A \otimes_K B$ : $(y \otimes v^{-1})(1 \otimes u) = y \otimes v^{-1}u = y \otimes (u-1)v^{-1}$, et $(y^{-1} \otimes v)(x \otimes 1 + 1 \otimes u) = (x-1)y^{-1} \otimes v + y^{-1} \otimes (u+1)v = (x \otimes 1 + 1 \otimes u)(y^{-1} \otimes v)$. L'exemple matriciel se vérifie par calcul direct, et $u_1 - 1 \neq u_1 + 1$ puisque $2 \neq 0$ dans $\mathbf F_3$.

f) La matrice de permutation est de trace nulle, et la trace sur l'espace nul est nulle ; en $s^2 = 1$ les caractères valent les dimensions $2$ et~$0$. Une fonction $P$ avec $P(1, 0) = 2$ et $P(1, 0) = 0$ n'existe pas.
\end{proof}

\begin{sloppypar}
\noindent\emph{Vérification en Lean.} a) est \texttt{finrank\_ne\_two}, pour tout corps algébriquement clos, puis \texttt{remark3\_fails} pour $\mathbf C$ ; b) est \texttt{two\_dvd\_mul}, \texttt{not\_isPrincipal\_I}, \texttt{T\_bijective}, \texttt{free\_I\_prod\_I} et \texttt{not\_free\_I} ; c) est \texttt{zero\_divisors} et \texttt{not\_isSimpleModule\_quaternion\_tensor} ; d) est \texttt{X\_ne\_pow} et \texttt{fixed\_not\_mem}, pour tout $p$ premier ; e) est \texttt{inv\_mul\_eq\_sub\_one\_mul}, \texttt{typed\_v1\_mul\_u1} et \texttt{corrected\_v1\_mul\_u1} dans tout $A \otimes_K B$, puis \texttt{presentation}, \texttt{prop12a\_typed\_fails} et \texttt{prop12a\_corrected}, vérifiés par \texttt{decide} ; f) est \texttt{character\_values} et \texttt{prop3\_fails}, dans \href{\leanblobmain/Redaction596Algebras.lean}{\texttt{Redaction596\allowbreak{}Algebras.\allowbreak{}lean}}. Ne sont pas formalisés : la structure de produit croisé (on vérifie seulement l'absence du sous-corps nécessaire), le module induit (on utilise la décomposition~(40) de la page), les puissances extérieures ($\chi_{\Lambda^0 E} = 1$ est pris tel quel), et le fait que le centre de $\mathbf H$ soit~$\mathbf R$.
\end{sloppypar}

\emph{Autres écarts relevés, non formalisés.} Au §~9, n\textsuperscript{o}\,10 (p.~162, et remarque~1 p.~163), l'ensemble des classes d'algèbres de quaternions est dit sous-groupe de $\mathrm{Br}(K)$ ; c'est faux en général, puisque le produit de deux classes de quaternions peut être la classe d'une algèbre à division de degré~$4$ (Albert) ; la construction est trop longue pour être reprise ici. Ces deux points ont été signalés par la lecture et ne sont pas revérifiés en détail. Aux n\textsuperscript{os}\,3 à~5 du §~9, la page prend pour $\Psi(\mathrm{cl}\,A[\mathcal E, L])$ la classe $[\mathcal E]$, alors que le bon signe est $[\mathcal E]^{-1}$ ; les théorèmes subsistent, l'inversion étant un automorphisme du groupe ; vérifier ce signe demanderait de formaliser les produits croisés. La démonstration du théorème~5 du §~12 (caractères de Brauer) n'a pas d'énoncé élémentaire faux à isoler, et Mathlib ne contient pas les caractères de Brauer.

\litterature{les énoncés justes sont classiques : Albert, \emph{Structure of Algebras} (1939) ; Serre, \emph{Corps locaux} (1962) et \emph{Représentations linéaires des groupes finis} (1967) ; pour $\mathbf Z[\sqrt{-5}]$, le théorème de Steinitz sur les modules projectifs de type fini sur un anneau de Dedekind. Le texte imprimé de 2012 n'a pas été consulté : on ne sait pas si ces écarts y sont corrigés.}
